
\subsubsection{2.1 Contamination Detection Methods}

\textbf{N-gram overlap detection.} Brown et al. [2020] introduced 13-gram exact match between benchmark items and training corpora as a contamination flag in the GPT-3 Appendix C. EleutherAI's lm-evaluation-harness \citep{Gaoetal2021} productionized this pipeline with an inverted index over sorted n-gram files. N-gram methods detect verbatim or near-verbatim contamination; paraphrased items evade detection \citep{Yangetal2023}.

\textbf{Embedding similarity detection.} Yang et al. [2023] showed that 8–18\% of HumanEval problems appear in RedPajama and StarCoder corpora in paraphrased form, detectable via embedding similarity. LLMSanitize \citep{Singhetal2024} provides a multi-method library integrating string-matching and embedding-based detection. Embedding methods require a meaningful similarity distribution over the corpus — an assumption that is examined in this paper.

\textbf{Membership inference attacks (MIAs).} Shi et al. [2024] introduced Min-K\% Prob, which identifies training members as sequences whose minimum-k\% token log-probabilities are high, without requiring corpus access. Zhang et al. [2024a] proposed Min-K\%++, an ICLR 2025 Spotlight that replaces a heuristic with a conditional categorical distribution mode criterion, achieving 6.2–10.5\% absolute AUROC improvement on WikiMIA. Zhang et al. [2024b] proposed DC-PDD, which computes cross-entropy divergence from a randomness reference. Fu et al. [2024], reviewing 50 papers on MIA-based detection, document conditions under which MIA methods perform at random guessing.

\subsubsection{2.2 Detector Inconsistency}

Singh et al. [2024] (ConTAM) evaluated 13 benchmarks across 7 models and found that different contamination metrics give inconsistent signals; the longest contaminated substring is identified as the most informative metric, without explaining when different metrics should be trusted. Dekoninck et al. [2024] (ConStat) frame contamination as non-generalizing performance. Xu et al. [2024] identify the absence of a unified comparison framework as an open problem in their comprehensive survey.

\subsubsection{2.3 Benchmark Contamination Audits}

Deng et al. [2024] found that GPT-4 achieves 57\% exact match on masked MMLU options, indicating memorization. Yang et al. [2023] showed 8–18\% of HumanEval appears in pretraining corpora in paraphrased form. Hidayat et al. [2025] conducted controlled leakage simulations on MMLU and HellaSwag, finding that n-gram detection achieves the highest F1 under controlled contamination.

\subsubsection{2.4 Positioning}

No prior work has proposed a geometry-based routing system for contamination detection, nor has prior work identified the stratum collapse phenomenon described in this paper.

